% !TeX program = xelatex
% ==========================================================================

% Пропоузал семестрового проєкту — курс «Штучний інтелект», УКУ.
% Це ЄДИНИЙ файл, який вам потрібно редагувати. Уся верстка й брендинг
% живуть у ucuproposal.cls — не змінюйте шрифти, поля, кеглі й кольори.
%
% Компіляція: XeLaTeX + Biber (Overleaf робить це автоматично).
% Локально:   latexmk (див. README.md).
% ==========================================================================
\documentclass[ukrainian]{ucuproposal}
% Опції класу (можна комбінувати через кому):
%   ukrainian | english  — мова документа й логотипу (за замовчуванням ukrainian)
%   slogan                — додає гасло УКУ в колонтитул титулки (за замовчуванням вимкнено)
\usepackage{hyperref}

% --------------------------------------------------------------------------
% Метадані титульної сторінки
% --------------------------------------------------------------------------
\projecttitle{Квантування, дистиляція та прунінг великих мовних моделей для української мови}
\projecttagline{Який метод стискання дає найкращий компроміс між якістю, розміром, пам’яттю та швидкістю – для українськомовної моделі?}

\teammember{Владислав Данилишин}
\teammember{Ія Магарита}
\teammember{Оксана Москвʼяк}
% \teammember[https://linkedin.com/in/...]{Прізвище Ім'я} % опційне клікабельне посилання на профіль (LinkedIn тощо)

\mentor{Марта Сумик | Олександр Косован}

\track{research} % research | engineering — не друкується на титулці, лише для внутрішньої категоризації

\repo{https://github.com/luftboud/small-LLM}

\proposaldate{2026-09-21}

\begin{document}
\maketitlepage

\section{Проблема та мотивація}
% Що саме ви розв'язуєте і чому це важливо й цікаво в контексті ШІ?
% Орієнтовно 1 абзац.
%Опишіть проблему, яку розв'язує ваш проєкт, та її значущість.
Сучасні LLM через велику кількість параметрів потребують значного обсягу пам’яті та обчислювальних ресурсів, що суттєво ускладнює їх локальне використання та розгортання на пристроях з обмеженими ресурсами.
Стиснення моделей за допомогою квантування, дистиляції та прунінгу може бути ефективним способом розв’язання цієї проблеми.
Метою проєкту є порівняти ці методи, їхні різновиди та комбінації, щоб визначити, які з них забезпечують найкращий баланс між розміром моделі, споживанням пам’яті, швидкістю інференсу та якістю роботи, зокрема з українською мовою.
\section{Огляд джерел та наявних рішень}
% Що вже зроблено у світі за цією темою? Спирайтесь на references.bib,
% цитуйте через \cite{}. Орієнтовно 3-5 джерел.
Стисло опишіть наявні підходи та роботи, на які ви спираєтесь~\cite{edge_centric_ai,edge_inteligence, optimizing_edge_ai, tiny_deep_learning}.



\section{Дані}
% Які набори даних, звідки, ліцензія. Якщо збираєте самі — як саме.
% Якщо дані не потрібні — поясніть чому.
\subsection{Тренування}
\textbf{Knowledge Distillation} потребує набору промптів, на яких велика модель дає цілі (текст або logits), а мала вчиться їх відтворювати. Датасет для цього від Lapa LLM на Hugging Face - \href{https://huggingface.co/datasets/lapa-llm/wiki-instruction-dialogs}{посилання тут}.

Структура датасету:
\begin{itemize}
    \item{Складений на основі україномовної Вікіпедії}
    \item{Складається з train частини (30000 записів), що містить такі поля}
    \begin{itemize}
        \item{task (тип задачі; string)}
        \item{instruction (формулювання завдання; string)}
        \item{input (фрагмент статті Вікіпедії, на якому виконується задача; string)} 
        \item{output (відповідь асистента; string)}
        \item{conversation (діалог з двох реплік; json)}
    \end{itemize}
\end{itemize}

\textbf{Pruning} та \textbf{Quantization} потребують невеликих даних для калібрації. Pruning потребує дані, аби зрозуміти, які частини моделі найменш значущі, щоб їх прибрати. Quantization потребує малих даних для того, аби визначити межі для нового типу даних, щоб зменшити похибку та зберегти ефективність моделі. Для цих задач можна взяти меншу частину датасету згаданого вище.

\subsection{Тестування}
Проєкт передбачає датасет для тестування точності мовної моделі до та після застосування оптимізаційних технік. Датасет для тестування взятий з Hugging Face, \href{https://huggingface.co/datasets/INSAIT-Institute/mmlu_ukr}{посилання тут}. 

Структура датасету:
\begin{itemize}
    \item{Складений на основі питань українського ЗНО та ЄВІ різних років}
    \item{Поділений на validation (1500 запитань) та test (14000 запитань) частини}
    \item{Кожна частина містить такі поля:}
    \begin{itemize}
        \item{question (зміст питання; string)}
        \item{subject (предмет, до якого стосується питання; string)}
        \item{choices (варіанти відповіді до запитання; list[string])} 
        \item{answer (правильна відповідь (індекс у списку choices); integer)}
    \end{itemize}
\end{itemize}

\section{Підхід і методи}
% Які алгоритми/архітектури плануєте використати, чи спираєтесь на
% існуючі реалізації та як плануєте їх покращити.
Для ефективної оптимізації мовної моделі, потрібно виконувати техніки у такому порядку: pruning, quantization, knowledge distillation.

\textbf{Pruning}.

\textbf{Quantization}.

\textbf{Knowledge Distillation}.


\section{Оцінювання результатів}
% Метрики, baseline, спосіб порівняння, якісні очікування.
\subsection{Baseline}
lapa-llm/lapa-v0.1.3-instruct 12B, text-only, BF16 як еталон якості.\cite{lapa_llm} (*спитати про початкове зменшення)
\subsection{Метрики і порівняння}
Традиційні моделі ML оцінюють переважно за точністю або F1-мірою. Моделі категорії TinyML вартує оцінювати за сукупністю метрик\cite{tiny_deep_learning}.
Ми плануємо оцінювати моделі з допомогою метрик якості й ефективності.

Основний перелік \textbf{метрик якості} виглядає так:
\begin{itemize}
    \item перевірка роботи моделі на різних датасетах/задачах;
    \item ROUGE-L\cite{rouge_l} - функція, яка порівнює summary моделі з еталоном XL-Sum.
\end{itemize}
А \textbf{метрики ефективності} це:
\begin{itemize}
    \item розмір моделі, GB;
    \item peak RAM/VRAM, GB;
    \item швидкість генерації, tokens/s;
    \item час від надсилання промпта до першого токену, ms
\end{itemize}

\textbf{good to have:} (це потім видалиться або переформатується)\\
Energy per generated token на RUBIK Pi\\
Blind human pairwise win rate

\subsection{Підхід до порівняння}
\begin{itemize}
    \item однакові prompts, decoding, splits, batch size, input/output lengths і тд.
    \item якість порівнювати по абсолютному score + retention
    \item ресурси порівнювати з baseline
    \item для benchmark scores виводити 95\% confidence intervals
    \item побудувати Pareto графіки + frontier
\end{itemize}
\section{План і ризики}
% Етапи до midterm і до фінального захисту, ключові ризики та план Б.
% Приклад врізки з ключовими цифрами (опційно, можна видалити):
\begin{keyfacts}
\textbf{Ключові цифри:} n датасети \textperiodcentered\ 13 тижнів \textperiodcentered\ ще щось
\end{keyfacts}

\begin{enumerate}
    \item \textbf{Підготовка}: Обрати baseline модель, провести аудит датасетів; Зафіксувати початкові показники базової моделі \textperiodcentered\ Тижні 1-2
    \item \textbf{Стискання}: Дослідити та застосувати INT8/INT4-квантизація, структурний прунінг і дистиляцію зі зменшенням кількості шарів \textperiodcentered\ Тижні 3-5
    \item \textbf{Порівняння та мідтерм}: Для кожної моделі отримати метрики; Мідтерм; Перевірити комбінації підходів та оцінити їх, визначити моделі з найкращим компромісом \textperiodcentered\ Тижні 6-8
    \item \textbf{Розгортання на edge-пристрої }: Конвертувати найкращу стиснену модель у сумісний формат і протестувати її запуск та продуктивність на платі RUBIС Pi 3; Отримати метрики \textperiodcentered\ Тижні 8-10
    \item \textbf{Аналіз}: Побудувати Pareto curves, провести human evaluation, error-analysis, забезпечити доказовість експерименту; Написати аналітичний звіт \textperiodcentered\ Тижні 9-10
    \item \textbf{Product phase}: Створити Comparison UI (як один промпт обробляється різними за розміром моделями); Упорядкувати репозиторій, написати вичерпну документацію; Створити постер та презентацію; \textit{Optional: написати англомовний preprint} \textperiodcentered\ Тижні 9-13
\end{enumerate}
\subsection{Ризики}
\begin{enumerate}
    \item 12B BF16 > 16 GB VRAM: CPU offload для baseline
    \item Два students і комбінації -- займає багато часу
    \item Немає фізичного локального Rubic Pi/NPU
\end{enumerate}
\section{Contribution statement}
% Хто за що відповідає. Відсотки не обов'язкові — якщо хочете їх вказати,
% додайте в квадратних дужках: \contribution[50\%]{Прізвище Ім'я}{...}
\begin{contributions}
  \contribution{Владислав Данилишин}{заглушка}
  \contribution{Ія Магарита}{заглушка}
  \contribution{Оксана Москвʼяк}{заглушка}
\end{contributions}
% варіанти: моделювання, евалюація, документація,дані, препроцесинг, EDA
% Примітка: декларація використання AI-інструментів (AI usage) до цього
% пропоузалу не входить — вона живе в README.md вашого репозиторію.

\ucuprintbibliography

\end{document}
